\section{The visible exact sector and its isolated spectrum}
\label{sec:exact}

The hidden-sector estimate leaves one possible source of slower decay:
exact gradients visible to fixed equilibrium calculations.  This section
constructs their deterministic operator, proves that nonstationary visible
projections lie in its weighted exact space, and isolates every mode below
the spin edge.  The vector space is deterministic; coefficients of its
vectors along a nonstationary trajectory may be random.  Throughout this section all horizons,
recipes, truncation levels, and approximation accuracies are fixed before
letting $n$ tend to infinity.  Successive approximation limits are then
taken in the order specified below.  The underlying disorder events are
those of \Cref{thm:path-calculus,thm:input-gap,prop:stat-posterior-gap}.
In particular, every use of an observation posterior gap refers to the
stopped, two-segment observation experiment of
\Cref{prop:stat-posterior-gap}.

\subsection{Two completions and a deterministic exact generator}

For stationary vector fields put
\begin{equation}\label{eq:exact-metrics}
 \ip{p}{q}_0=\E_\mu p^{\mathsf T}q,
 \qquad
 \ip{p}{q}_v=\E_\mu p^{\mathsf T}D_vq,
 \qquad
 \mathfrak d(p)=\E_\mu\sum_{i,j}v_j\abs{d_jp_i}^2.
\end{equation}
We use the same notation for the deterministic limits of these forms on
recipes.  The first metric defines the visible space already constructed;
completing it in the weaker second metric gives $\mathcal H_v$.  The map $\iota:\mathcal H\longrightarrow\mathcal H_v$ is
continuous and injective.  To see injectivity, multiply first by a smooth
cutoff supported on $v_i\ge\varepsilon$.  On that support the unweighted
norm is at most $\varepsilon^{-1/2}$ times the weighted norm.  For a fixed
recipe the removed unweighted mass tends to zero as $\varepsilon\downarrow0$,
by the spatial square-tail estimates; bounded multipliers and dense
approximation extend this assertion to each fixed element of $\mathcal H$.
Thus equality in the weighted completion identifies ordinary visible
vectors, but weighted approximation need not imply ordinary approximation.
We will use the forward embedding only; its inverse need not be bounded.

Write $L^{(n)}$ also for the componentwise scalar generator on vector
fields, and write $\mathcal G^{(n)}$ for the exact gradient extension in
\Cref{eq:path-generator}.  The same letters without the superscript will
be used for the limiting operators when their domain is unambiguous.

\begin{lemma}[Semigroups and smoothing]\label{lem:exact-smoothing}
The stationary limits define strongly continuous semigroups $K_t$ on
$\mathcal H$ and on $\mathcal H_v$, and
$\iota K_t=K_t\iota$.  For each fixed $T$ their operator norms on both
spaces are bounded uniformly for $0\le t\le T$.  On $\mathcal H$ the
semigroup is analytic and, for $t>0$,
\begin{equation}\label{eq:exact-graph}
 \norm{L K_t}_{\mathcal H\to\mathcal H}
 +\norm{(-L)^{1/2}K_t}_{\mathcal H\to\mathcal H}
 \le C_t.
\end{equation}
The corresponding finite-dimensional bounds are uniform in $n$ on the
retained disorder events.  Here $L$ is the generator of componentwise
stationary heat flow, rather than the generator of $K$.
\end{lemma}

\begin{proof}
We first establish local boundedness in each norm, and then smoothing in
the ordinary norm.  Conditional prediction and the path expansion in
\Cref{thm:path-calculus} preserve the visible space.  In the ordinary norm,
\Cref{prop:path-rough} gives the bound and removes likelihood tails in
stationary square mean.  For the weighted bound, integrate the square
evolution against $\mu$ and retain the weights at both edge corners.
The principal integration-by-parts term absorbs a fixed multiple of
\[
 \E_\mu\sum_{i,j}v_iv_j\abs{d_jp_i}^2.
\]
Here both corners of an edge are retained.  The elementary bounds
\[
 v_i(x^j)/v_i(x)\in[e^{-C|J_{ij}|},e^{C|J_{ij}|}],
 \qquad |d_jv_i(x)|\le C|J_{ij}|v_i(x),
 \qquad \sup_i\sum_jJ_{ij}^2\le C
\]
control the terms differentiating the weight.  Cauchy--Schwarz and Young's
inequality absorb the first-order term using
$\sum_j a_{ij}^2/v_j\le C$.  The multiplication term is bounded by
$\norm{aD_v^{-1}}\le C$, the comparable edge weights, and $v\le1$.
Consequently
\[
 \frac{\dd}{\dd t}\norm{K_tp}_v^2
 \le C\norm{K_tp}_v^2
   -c\E_\mu\sum_{i,j}v_iv_j|d_j(K_tp)_i|^2.
\]
This proves the extension to the weighted completion.  Strong continuity
holds on smooth recipes by their difference bounds and then on both
completions by density and local boundedness.

Analyticity follows by treating the first-order and multiplication terms
as a relatively bounded perturbation of componentwise heat flow.  Their
uniform relative bound is
\begin{equation}\label{eq:exact-relative}
 \norm{(\mathcal G^{(n)}-L^{(n)})p}_0
 \le C\bigl(\norm{(-L^{(n)})^{1/2}p}_0+\norm{p}_0\bigr).
\end{equation}
For example, square the first-order term row by row and use
$\sum_j a_{ij}^2/v_j\le C$; its sum is bounded by $C\mathfrak d(p)$.
The remaining term has bounded operator norm.  If $z$ is on either ray
bounding a sector strictly larger than the negative real half-line, the
spectral theorem gives
\[
 \norm{(-L^{(n)})^{1/2}(z-L^{(n)})^{-1}}
 \le C|z|^{-1/2},\qquad
 \norm{(z-L^{(n)})^{-1}}\le C|z|^{-1}.
\]
At sufficiently large $|z|$, independently of $n$, the resolvent for
$\mathcal G^{(n)}$ is therefore given by a convergent Neumann series.
Contour inversion gives uniform analytic bounds.  Moreover,
\Cref{eq:exact-relative} and the scalar inequality
$\sqrt u\le\varepsilon u+C_\varepsilon$ show that the graph norms of
$L^{(n)}$ and $\mathcal G^{(n)}$ are uniformly equivalent after adding
a fixed multiple of the identity.  This proves \Cref{eq:exact-graph}.
Passing first on the dense recipe class and then by the uniform bounds
gives the assertions in the limiting spaces.
\end{proof}

We now choose exact test gradients whose potentials will also detect
every weighted-visible exact limit.  For a smooth bounded scaled vector
recipe $B$ and $\delta>0$, define
\begin{equation}\label{eq:exact-potential}
 f_{B,\delta}
 =\delta^{-1}\left\{x^{\mathsf T}S_\delta B
       -S_\delta\bigl[(S_\delta x)^{\mathsf T}B\bigr]\right\}.
\end{equation}
The semigroup in the first term acts componentwise.  For every scalar
test $G$, reversibility gives the pairing
\[
 \ip{G}{f_{B,\delta}}_{L^2(\mu)}
 =\delta^{-1}\E_\mu\!\left[
 B^{\mathsf T}\{S_\delta(xG)-(S_\delta x)S_\delta G\}\right].
\]
The expression in braces is the conditional covariance of the spin
vector and $G$ after a time $\delta$, given the initial state.  Thus
$f_{B,\delta}$ tests short-time covariance in the direction $B$;
in particular it is centered.  For fixed $B,\delta$, the potentials have
polynomial supremum bounds and
\begin{equation}\label{eq:exact-test-bound}
 \sup_x\norm{df_{B,\delta}(x)}\le C_{B,\delta}.
\end{equation}
Indeed, the derivative matrix of $B$ has operator norm $O(n^{-1/2})$
and $\norm{x}=\sqrt n$.  The required derivatives of $S_\delta B$ and
$S_\delta x$ are direct single-future derivatives: for each fixed output
combination $c$, use $dS_\delta(c^{\mathsf T}F)
=K_\delta d(c^{\mathsf T}F)$ and the point-start bound in
\Cref{prop:path-rough}.  Taking the supremum over $c$ then bounds the
already averaged derivative matrix at each starting state.  The exact
flip product rule controls its quadratic term by the column square
bounds.  Apply the same single-future propagation to the scalar output
$(S_\delta x)^{\mathsf T}B$ to obtain \Cref{eq:exact-test-bound}.

To put these derivatives in the visible space, first truncate the
conditionally averaged derivative matrices by common linear path
truncations, and only then open their future branches.  This is the single-future case of
\Cref{lem:path-conditional-derivative}; it requires no inference from
separate slot Jacobians of a general multitime function.  Products and
diagonals of the resulting bounded path matrices are covered by
\Cref{thm:path-calculus}.  Whenever a fixed smooth multitime circuit is
used in a further conditional derivative, its mixed-slot bounds and
common tensor truncation are those of
\Cref{lem:path-mixed-tensors,lem:path-conditional-derivative}, and its
opened matrix belongs to the path closure by
\Cref{prop:path-opened-derivative}.  All such circuit choices precede
dimension.  Thus $df_{B,\delta}$ and its $K$-orbits are well-defined
deterministic vectors of $\mathcal H$.

\begin{definition}\label{def:exact-spaces}
Let $\mathcal X$ be the closure in $\mathcal H$ of the linear span of
\[
 K_tdf_{B,\delta},\qquad t\ge0,\quad\delta>0,
 \quad B\text{ a smooth bounded scaled recipe}.
\]
Let $\mathcal X_v$ be the closure of the same span in $\mathcal H_v$.
Equivalently, countable dense recipe, time, and accuracy sets may be used.
\end{definition}

\begin{theorem}[The exact visible operator]\label{thm:exact-sector}
The spaces in \Cref{def:exact-spaces} and their semigroups are deterministic.
The embedding $\iota:\mathcal X\to\mathcal X_v$ is continuous, injective,
and dense.  On $\mathcal X_v$, $K_t$ is a positive self-adjoint semigroup.
If $\mathcal A_v$ is minus its generator, then
$\spec(\mathcal A_v)\subset[\gamma,\infty)$ for the fixed gap constant
$\gamma>0$ in \Cref{thm:input-gap}.  In particular the number
\begin{equation}\label{eq:exact-lambda}
 \lambda
 =\min\left\{\lambda_{\rm sp},\inf\spec(\mathcal A_v)\right\}
 \quad\text{satisfies}\quad
 0<\lambda\le\lambda_{\rm sp}\le(1-\beta)^2.
\end{equation}
All dimension limits in this definition precede the spectral infimum.
\end{theorem}

\begin{proof}
On the generating span, the path limits make every inner product and
fixed-time matrix element deterministic.  Self-adjointness comes from
finite-dimensional exactness, through the identity
\[
 \ip{df}{K_tdg}_v
 =\ip{f}{(-L)S_tg}_{L^2(\mu)}.
\]
The right side is symmetric in $f,g$, is nonnegative when $f=g$, and
has operator norm at most $e^{-\gamma t}$ in the gradient norm.
These properties pass to the weighted closure.  The spectral theorem
then gives the generator and its lower bound.  The embedding properties
follow from \Cref{eq:exact-metrics} and the definition of the closures;
the final spin-edge bound is \Cref{prop:spin-spectrum}.
\end{proof}

\subsection{Saturation by stationary exact gradients}

\begin{lemma}[Saturation]\label{lem:exact-saturation}
If a weighted-visible vector is a weighted mean-square limit of stationary
exact gradients, then it belongs to $\mathcal X_v$.  In particular this
holds for any such vector initially given in $\mathcal H$.
\end{lemma}

\begin{proof}
It suffices to show that the component orthogonal to $\mathcal X_v$
vanishes.  Subtract the weighted orthogonal projection and call the
remainder $q$.  It remains weighted-visible and a weighted limit of exact
gradients, while its pairing with every generating orbit is zero.  Take centered potentials $G_n$
representing this exact limit.  The gap gives
$\norm{G_n}_{L^2(\mu)}^2\le\gamma^{-1}\norm{dG_n}_v^2$.
For every fixed test potential $f=f_{B,\delta}$,
\begin{equation}\label{eq:exact-potential-pairing}
 \ip{G_n}{f}_{L^2(\mu)}
 =\int_0^\infty\ip{dG_n}{K_tdf}_v\dd t.
\end{equation}
The exponential gap bound gives a uniform integrable tail; the limit is
zero by orthogonality.

The same potential pairing has a local form that detects $q$ as the
short time interval shrinks.  Reversibility and conditional martingale
covariance give the exact identity
\begin{align}\label{eq:exact-covariance-potential}
 \ip{G_n}{f_{B,\delta}}_{L^2(\mu)}
 =\frac2\delta\int_0^\delta
 \E_\mu\!\left[
 B(X_0)^{\mathsf T}(dS_{\delta-s}x)(X_s)
 D_{w(X_s)}dS_{\delta-s}G_n(X_s)\right]\dd s.
\end{align}
Here the derivative matrix has output as row.  The factor $2$ is the
quadratic-variation factor for flip rate $w_j/2$ and spin increment $2$.
The derivative matrix $dS_rx$ differs from the identity by $O(r)$ in
operator norm: integrate $dS_rx-I=\int_0^r dS_uLx\dd u$ and use
$\norm{dLx}\le C$ together with the pointwise rough gradient bound.
Also $\E_\mu\norm{B(X_s)-B(X_0)}^2=O_B(s)$.

The passage $\delta\downarrow0$ is justified in the weighted norm,
without comparing the two completions in reverse.  For every own-spin
independent vector $p$,
\begin{equation}\label{eq:exact-w-isometry}
 \E_\mu\norm{D_wp}^2
 =\E_\mu\sum_i\E_\mu[w_i^2\mid x_{-i}]p_i^2
 =\E_\mu\sum_i v_ip_i^2,
 \qquad
 \E_\mu[w_i\mid x_{-i}]=v_i.
\end{equation}
Thus $D_w$ is an isometry from the closure of exact gradients with their
weighted norm into ordinary square mean.  In \Cref{eq:exact-covariance-potential}
one may first replace $G_n$ by $S_\varepsilon G_n$, shrink $\delta$, and
then shrink $\varepsilon$, using strong continuity on exact gradients.
Equivalently one passes in the displayed covariance using
\Cref{eq:exact-w-isometry}.  It follows from the zero pairing that
$\ip{B}{D_wq}_0=0$ for every $B$.

To conclude from these zero pairings, we must place $D_wq$ in the ordinary
visible space.  Cut off the coordinates where $v_i<\varepsilon$.  On the retained set a weighted-visible approximation
is an ordinary visible approximation, and multiplication by $w$ is an
allowed bounded label.  On the removed set \Cref{eq:exact-w-isometry}
and own-spin independence identify its squared norm with the corresponding
weighted mass of $q$.  That mass tends to zero for a fixed element of
$\mathcal H_v$, by density of bounded recipes.  Hence $D_wq\in\mathcal H$.
Its orthogonality to all recipes makes it zero, and the isometry makes
$q=0$.
\end{proof}

\subsection{Random coefficients and a reverse entropy budget}

The stationary construction now has the required saturation property:
a weighted-visible limit of exact gradients lies in $\mathcal X_v$.
The remaining task is to prove that the visible projection of an
\emph{evolving} gradient is such a limit.  Its coefficients can be random,
so stationary exactness cannot simply be projected onto each recipe.
We first retain those coefficients in a scalar Hilbert space.  A reverse
entropy estimate will then let us test their vector directions for
weighted curl; the observation inequality below converts vanishing curl
into approximation by gradients.

Consider an ordinary trajectory on a fixed interval $[a,b]$ whose initial
law has relative log density $o(n)$ with respect to $\mu$, in the
finite-tree convention of \Cref{thm:path-calculus}.  Let
\[
 p_s=dS_{D-s}f,\qquad a\le s\le b\le D,
 \qquad \E\norm{p_b}^2\le C,
\]
with polynomial pointwise bounds throughout.  At every time used in the
argument, retain
the countable family of bounded-$L^2$ scalar sequences consisting of
recipe contractions with $p_s$ and the scalar tests to be used below.
Close these families under rational linear combinations and the actual
conditional expectations between the retained times.  Extract a
subsequence on which all pairwise scalar inner products converge,
quotient by sequences of limiting norm zero, and complete.  This defines
a Hilbert space $\mathcal S_s$ at time $s$.  Conditional expectation
descends to a contraction
$C_{s,t}:\mathcal S_t\to\mathcal S_s$, because its finite-dimensional
$L^2$ norm is at most one.  Additional scalar tests can be included
before extraction; they need not be regular functions of the spins.

The extraction is joint for any prescribed finite or countable family
of fixed comparison times; that family can be enlarged before extraction.
We also retain the comparison-time sequences selected below from the
sets where the integrated typicality tests and norm budgets hold.
After a further subsequence their times converge.  Their scalar spaces
are defined from the tests and conditional expectations at the actual
finite-dimensional times, with separate copies for distinct sequences
even if their limiting times coincide.  We do not select arbitrary
dimension-dependent times from exceptional sets or identify these spaces
as a previously constructed continuous-time process.  At positive limiting
time separation, strong continuity and fixed smoothing identify the
vector semigroup with $K$ at the limiting separation.  The selection is
made in \Cref{lem:exact-good-times} and in the proof of
\Cref{prop:cutoff-block}.

Here is the continuity estimate that permits rough random coefficients.
For fixed bounded scaled recipes $B,B'$, typicality gives concentration
of $\norm{B-B'}^2$ at its deterministic Gram limit.  On that event,
Cauchy--Schwarz bounds the squared contraction by
$\norm{p_s}^2\norm{B-B'}^2$.  The polynomial pointwise bounds and the
upper relative log-density $o(n)$ remove the exponentially small
exceptional set.  Consequently
\begin{equation}\label{eq:exact-contraction-continuity}
 \limsup_n\E|p_s^{\mathsf T}(B-B')|^2
 \le\norm{B-B'}_{\mathcal H}^2
                   \limsup_n\E\norm{p_s}^2.
\end{equation}
Thus contraction extends continuously from recipes to their completion;
this conclusion uses nonstationary typicality, not only stationary
mean-square approximation.

Finite projections on an orthonormal recipe basis now have coefficients
in $\mathcal S_s$ given by their pointwise contractions with $p_s$.
Typical Gram convergence and Bessel's inequality bound the sum of their
scalar second moments.  Increasing the finite basis therefore defines
\[
 q_s\in\mathcal H\mathbin{\widehat\otimes}\mathcal S_s,
 \qquad \norm{q_s}^2\le\liminf_n\E\norm{p_s}^2.
\]
The first tensor factor records the equilibrium recipe direction; the
second records its possibly rough random coefficient.  Conditional
expectation will act on the latter, while all differentiation acts on the
former.

\begin{proposition}[Transport of the visible coefficients]
\label{prop:exact-transport}
For prescribed comparison times $a\le s<t\le b$ retained in the
preceding construction,
\begin{equation}\label{eq:exact-transport}
 q_s=(K_{t-s}\mathbin{\widehat\otimes}C_{s,t})q_t.
\end{equation}
The identity also holds for the retained selected-time sequences with
positive limiting separation: $K_{t-s}$ uses that limiting separation,
while $C_{s,t}$ is induced by their actual conditional expectations.
For an interior time $s$ with a fixed positive amount of future evolution,
$q_s$ belongs to the graph domain of componentwise $L$ in the vector
factor, with graph norm bounded by $C_{t-s}\norm{q_t}$.  No derivative is
taken on the scalar factor.
\end{proposition}

\begin{proof}
We test the claimed transport first on finite recipe projections.
Truncate the path expansion of $K_{t-s}$ at fixed order and likelihood
level.  Between finitely many testing recipes, its matrix elements
concentrate at their deterministic stationary values by
\Cref{thm:path-calculus}.  To remove the endpoint complement, apply the
adjoint bounded path word to a testing recipe and project at time $t$.
Conditional prediction on replicated trees makes the remaining squared
contraction arbitrarily small.  This squared estimate is obtained
\emph{before} multiplying by any scalar test; Cauchy--Schwarz therefore
permits arbitrary bounded-$L^2$ scalar coefficients, including
past-dependent tests.  Polynomial pointwise bounds and the exponential
exceptional-set estimates justify discards.  The operator-small path
tails then remove the truncations.  Increasing the finite basis gives
\Cref{eq:exact-transport}, since conditional expectation is a contraction.
Apply \Cref{eq:exact-graph} to the first tensor factor for the last
assertion.
\end{proof}

\begin{lemma}[Good comparison times]\label{lem:exact-good-times}
Let $h_s$ be the marginal density relative to $\mu$, and set
$r_{s,j}(x)=h_s(x^j)/h_s(x)$.
In every fixed open subinterval $I\Subset(a,b)$ one can choose comparison
times $s=s_n\in I$, after extraction, for which
\begin{align}
 &\E\norm{p_s}^2+
   \E\sum_jw_j\norm{d_jp_s}^2\le C_I,
 \label{eq:exact-hessian-budget}\\
 &\E\norm{p_s}^2\sum_jw_j(1-\sqrt{r_{s,j}})^2=o(n).
 \label{eq:exact-reverse-budget}
\end{align}
All fixed typicality tests and their bounded auxiliary-time integrals
can be retained at these times.
\end{lemma}

\begin{proof}
We obtain the two budgets from different estimates on the same path.
The rough square estimate supplies the time-integrated first line and
$\norm{p_s(x)}^2\le C\E[\norm{p_b(X_b)}^2\mid X_s=x]$.
Tilt the entire path by
$g(X_b)/\E g(X_b)$, where $g=1+\norm{p_b}^2$.
Its entropy relative to the stationary path is $o(n)$: the original
relative log density is $o(n)$ and the extra log density is $O(\log n)$.
The tilt depends on the forward terminal state, so it preserves the
original reverse transitions.  Those transitions have rates
$w_jr_{s,j}/2$, whereas stationary reverse rates are $w_j/2$.
The jump-process entropy chain rule and
\[
 (\sqrt r-1)^2\le r\log r-r+1\qquad(r\ge0)
\]
give an integrated Hellinger budget $o(n)$ under the tilted path.
The preceding conditional bound on $p_s$ converts it into the time
integral of \Cref{eq:exact-reverse-budget}.  Markov's inequality selects
a positive-measure set of times where the Hessian budget is bounded and
the reverse budget is $o(n)$.  Countably many fixed typicality tests can
be imposed successively there by their integrated failure bounds and
extraction.  Auxiliary-time exceptional sets cost their square-mean
bounds, rather than a polynomial supremum times their measure.
\end{proof}

\subsection{Detecting exactness without differentiating coefficients}

The next argument tests curl through stationary adjoints.  Its error
under a nonstationary law is controlled by the reverse entropy budget.
To retain that error explicitly, for a law $h\mu$ put
$r_j(x)=h(x^j)/h(x)$.  Pairing the two corners of
each $j$-edge gives, for arbitrary $F_j$,
\begin{equation}\label{eq:exact-nonstationary-ibp}
 \E_{h\mu}\sum_j\bigl[(x_j-m_j)F_j-v_jd_jF_j\bigr]
 =\E_{h\mu}\sum_j
 (x_j-m_j)\frac{w_j}{2}(1-r_j)F_j.
\end{equation}
The formula extends by continuity if $h$ vanishes.  At positive times
all its ratios are finite.

Let $A$ be a bounded smooth antisymmetric word matrix from the matrix
calculus, with its column-flip versions retained in that calculus.  Set
\begin{align}\label{eq:exact-adjoint}
 (B_A)_i
 =\frac1{\sqrt n}\sum_j\left[
 (x_j-m_j)(v_iA_{ij})(x^j)
 -v_jd_j\bigl((v_iA_{ij})(x^j)\bigr)\right].
\end{align}
The adjoint identity at stationarity is
\begin{equation}\label{eq:exact-curl-adjoint}
 \ip{q}{B_A}_0
 =\E_\mu\sum_{i,j}\frac{v_iv_j}{\sqrt n}A_{ij}d_jq_i.
\end{equation}
Choose $A$ first as a finite one-state smooth word test.  Then $B_A$ is
visible: its first term is a matrix word applied to a bounded site vector
with the stated normalization, and the tied-index derivatives in the
second term are the one-state column-flip and Hadamard operations of
\Cref{thm:path-calculus}.  Conditional matrix limits are reached only
after these fixed tests, through the strong approximation below.  If a
conditional derivative is opened in forming a fixed test approximant,
\Cref{prop:path-opened-derivative} supplies its path-matrix closure from
the mixed-slot circuit construction; no general first-Jacobian criterion
is being invoked.

\begin{lemma}[Annihilation of the stationary adjoints]
\label{lem:exact-adjoints}
At a comparison time from \Cref{lem:exact-good-times}, every scalar
component of $q_s$ annihilates $B_A$.  It also annihilates the stationary
adjoints of
$\E_\mu\sum_i v_i b_i d_iq_i/\sqrt n$, for bounded smooth diagonal labels
$b$.  These are scalar-square-mean assertions.
\end{lemma}

\begin{proof}
The conclusion concerns the square mean of each random coefficient, so
we multiply an adjoint test by a second visible contraction.  Write
$p=p_s$, and take $e=p^{\mathsf T}B$ with a fixed smooth scaled recipe $B$.  The flip product rule and
\Cref{eq:exact-hessian-budget} give
$\E\sum_jw_j|d_je|^2\le C_B$.
In \Cref{eq:exact-nonstationary-ibp} use
\[
 F_j=e\,U_j,\qquad
 U_j=\frac1{\sqrt n}\sum_i p_i(v_iA_{ij})(x^j).
\]
The elementary column estimate needed here is
\begin{equation}\label{eq:exact-column-budget}
 \E\sum_jw_j(1+\sqrt{r_j})^2|U_j|^2
 \le\frac{C_A}{n}
 \left(\E\norm{p}^2+\E\sum_jw_j\norm{d_jp}^2\right).
\end{equation}
For the term without $r_j$, this is the bounded operator estimate for
the column-flipped test matrix.  For the term with $r_j$, change $x$ to
$x^j$ using detailed balance.  Replace
$p_i(x^j)$ by $p_i(x)-2x_jd_jp_i(x)$; the first part has the same
operator bound, and bounded individual column norms bound the second
part by the displayed derivative budget.  Local comparability of $v$
at the two corners supplies the same constants for the weighted labels.

Factor $1-r_j=(1-\sqrt{r_j})(1+\sqrt{r_j})$ in the remainder.
Since $|e|\le C_B\norm{p}$ and $|x_j-m_j|\le2$,
\Cref{eq:exact-reverse-budget,eq:exact-column-budget} make that remainder
$o(1)$ by Cauchy--Schwarz: its two squared budgets are $o(n)$ and
$O(n^{-1})$.  Derivatives falling on $e$ cost $O(n^{-1/2})$, by its
Dirichlet bound and the same column estimate.  Derivatives falling on
$p$ cancel exactly, since $d_jp_i=d_ip_j$ and $A$ is antisymmetric.
The remaining identity is
\[
 \E\bigl[(p^{\mathsf T}B)(p^{\mathsf T}B_A)\bigr]=o(1).
\]
Approximate $B_A$ in the visible norm by allowed recipes $B$ and use
Bessel's inequality and the squared contraction comparison before
removing that approximation.  The result is
$\norm{\ip{q_s}{B_A}_{\mathcal H}}_{\mathcal S_s}^2=0$.
The diagonal argument is identical, with the cancellation supplied by
$d_ip_i=0$.  It uses the stationary adjoint of the stated diagonal
functional and the same $n^{-1/2}$ normalization.
\end{proof}

Adjoint annihilation is initially a collection of weak tests.  To turn
it into a vanishing curl norm, we next approximate the derivative of the
\emph{limiting visible vector} strongly.  This excludes derivative mass
that escapes every fixed word test.

\begin{lemma}[Strong derivative closure]\label{lem:exact-form-closure}
Let $q$ be a deterministic vector component of an interior $q_s$, obtained
by contraction with any fixed scalar Hilbert direction.  Then the pair
\[
 \left((\sqrt{v_iv_j}\,d_jq_i)_{ij},
       (\sqrt{v_i}\,d_iq_i)_i\right)
\]
is a strong mean-square limit of normalized smooth word matrices and
their diagonal parts.  Consequently the adjoint annihilation in
\Cref{lem:exact-adjoints} implies that representatives $Q^{(n)}$ of $q$
can be chosen with bounded $\mathfrak d(Q^{(n)})$ and
\begin{equation}\label{eq:exact-vanishing-curl}
 \E_\mu\sum_{i,j}v_iv_j|d_jQ_i^{(n)}-d_iQ_j^{(n)}|^2
 +\E_\mu\sum_i v_i|d_iQ_i^{(n)}|^2\longrightarrow0.
\end{equation}
Each successive finite approximant may be chosen polynomially bounded.
\end{lemma}

\begin{proof}
The approximation has two stages: smoothing in form norm, followed by
finite recipe approximation.  By
\Cref{prop:exact-transport,lem:exact-smoothing}, $q$ lies in the graph
domain of componentwise $L$.  First replace it by $S_\varepsilon q$.
The spectral theorem makes this approximation strong in the norm
$\norm{q}_0^2+\mathfrak d(q)$.  For fixed $\varepsilon>0$, approximate the
input $q$ by finite smooth recipes $q_k$.  Since
\[
 \norm{(-L)^{1/2}S_\varepsilon(q-q_k)}_0
 \le C_\varepsilon\norm{q-q_k}_0,
\]
the second approximation is also strong in form norm.  This norm
controls both weighted derivative quantities displayed in the statement.
These operations apply to the deterministic vector $q$; the scalar
Hilbert direction used to obtain it has not been differentiated.

For fixed $\varepsilon,k$, the derivative of the one-state recipe
$q_k$ is $A/\sqrt n$ in the smooth word, Hadamard, and column-flip
closure.  Bounded finite-rank errors disappear in normalized
Hilbert--Schmidt norm.  Propagate this terminal derivative by the direct
single-future identity $dS_\varepsilon=K_\varepsilon d$, separately for
each component.  The common columnwise path tail bound approximates this
operation in Hilbert--Schmidt square mean, because the fixed input
derivative has bounded Hilbert--Schmidt norm.

The remaining conditional expected path matrices lie in the same strong
tracial closure.  Choose their bounded labels as fixed smooth multitime
circuits in the class of \Cref{lem:path-mixed-tensors}, with likelihood
caps and any finite spin-sign decomposition already fixed.  Then
simulate the finitely many transitions with the nested accuracies of
\Cref{thm:path-calculus}.  Conditional pair-trace tests give the squared
error.  Averaging $r$ independent simulation seeds has conditional
Monte Carlo error at most $C/r$ in normalized Hilbert--Schmidt square
mean.  Each fixed simulated calculation is a one-state word matrix.
Only strong Hilbert--Schmidt approximation is required in this step:
we neither differentiate a spatial-square-mean approximation of a label
nor promote its error to an operator error.  When an actual derivative
of a fixed multitime circuit is needed elsewhere, its operator
truncation and subsequent matrix reduction are instead the separate
conclusions of
\Cref{lem:path-conditional-derivative,prop:path-opened-derivative}.
Thus the derivative approximation here is strong, rather than only a
list of converging individual traces.  Diagonal extraction is performed
before taking closure and is a contraction for the relevant diagonal
norm; its convergence also follows directly from the form approximation.

The derivative matrices were strongly approximated by $A/\sqrt n$
with $A$ in the word closure.  Multiplication on the two sides by
$D_{\sqrt v}$ is a contraction in Hilbert--Schmidt norm.  The weighted
curl is therefore strongly approximated by
$D_{\sqrt v}(A-A^{\mathsf T})D_{\sqrt v}/\sqrt n$.
Since the word closure is linear and transpose-stable, these are
precisely weighted antisymmetric test matrices of the type used in
\Cref{eq:exact-curl-adjoint}.
\Cref{eq:exact-curl-adjoint} and adjoint annihilation force its inner
product with every element of this dense antisymmetric test class to be
zero, hence its norm is zero.  The diagonal tests do the same for the
own-spin derivative.  The strong form approximants then give
\Cref{eq:exact-vanishing-curl}.  Applying the argument to all scalar
Hilbert directions, or to a countable orthonormal basis of their
separable span, proves the corresponding tensor assertion in square
mean.
\end{proof}

\begin{lemma}[Removing own-spin dependence]\label{lem:exact-own-spin}
Let $\overline Q_i=P_iQ_i=\E_\mu[Q_i\mid x_{-i}]$, and put
\[
 O(Q)=\E_\mu\sum_i v_i|d_iQ_i|^2,\qquad
 C(Q)=\E_\mu\sum_{i,j}v_iv_j|d_jQ_i-d_iQ_j|^2.
\]
Writing $\delta_J=\max_{i,j}|J_{ij}|$ and
$R_J=\max_i\sum_jJ_{ij}^2$, one has
\begin{align}\label{eq:exact-own-spin-bound}
 \norm{Q-\overline Q}_v^2&\le O(Q),\notag\\
 C(\overline Q)&\le3e^{12\delta_J}C(Q)
                    +12e^{14\delta_J}R_JO(Q).
\end{align}
Thus the representatives in \Cref{lem:exact-form-closure} may be made
own-spin independent without changing their weighted limit or the
vanishing-curl conclusion.
\end{lemma}

\begin{proof}
The first assertion is the conditional-variance identity
\[
 \norm{Q-\overline Q}_v^2
 =\E_\mu\sum_i v_i^2|d_iQ_i|^2\le O(Q).
\]
For the curl estimate, work on one conditional two-spin law before
summing over pairs.  Condition on the spins outside $i,j$ and write the
pair law as $p_{ab}$, $a,b\in\{-1,1\}$, with coupling $K=J_{ij}$.
Delete $K$ and normalize to obtain a product law $q=q_i\otimes q_j$.
Put $m=\E_q a$, $m'=\E_q b$, $A=1-m^2$, and $B=1-(m')^2$.
Uniformly in the fields,
\[
 \frac{p_{ab}}{q_{ab}},\quad\frac{v_i(b)}A,
 \quad\frac{v_j(a)}B\ \in[e^{-2|K|},e^{2|K|}].
\]
Write $u_b=d_iQ_i(b)$ and $t_a=d_jQ_j(a)$.  Affinity in one spin gives
\[
 (P_iQ_i)(b)=\E_{q_i}Q_i(a,b)+(m_i(b)-m)u_b.
\]
Consequently the averaged curl has the exact representation
\begin{equation}\label{eq:exact-average-difference}
 d_j\overline Q_i-d_i\overline Q_j=C_0+E_i-E_j,
 \qquad C_0=\E_q[d_jQ_i-d_iQ_j],
\end{equation}
where
\[
 E_i=\tfrac12\{(m_i(+)-m)u_+-(m_i(-)-m)u_-\}
\]
and $E_j$ has the symmetric definition.  In particular all three
quantities in \Cref{eq:exact-average-difference} are independent of the
two spins.  With $W=\E_p[v_i(b)v_j(a)]$, comparison and Jensen give
\[
 WC_0^2\le e^{12|K|}\E_p[v_iv_j|d_jQ_i-d_iQ_j|^2].
\]
The log derivative of $\operatorname{sech}^2$ has absolute value at most
$2$, so $|m_i(\pm)-m|\le|K|e^{2|K|}A$.  Thus
\[
 E_i^2\le\tfrac12K^2e^{4|K|}A^2(u_+^2+u_-^2),
 \qquad B(u_+^2+u_-^2)\le4\E_{q_j}u_b^2.
\]
The last inequality retains the two opposite corners even when one has
very small probability.  Combining it with $W\le e^{6|K|}AB$,
$A^3\le A$, and the comparison of $p\,v_i$ with $qA$, gives
\[
 WE_i^2\le2K^2e^{14|K|}\E_p[v_i|d_iQ_i|^2].
\]
Use the analogous bound for $E_j$, square
\Cref{eq:exact-average-difference} with factor $3$, average over the
other spins, and sum ordered pairs.  The two error sums together cost
$12e^{14\delta_J}R_JO(Q)$, proving the claim.  No lower bound on an
individual conditional variance, and no unweighted derivative bound,
is used in this removal step.
\end{proof}

\subsection{An observation inequality excluding a harmonic defect}

We now need a uniform estimate converting small weighted curl into
small weighted distance from gradients.  The observation process provides
that estimate: at its product endpoint a Fourier calculation gives it
directly, and the posterior gap controls its transport back to the
original law.

\begin{lemma}[No harmonic defect]\label{lem:exact-no-defect}
Let $Q=Q^{(n)}$ be a polynomially bounded vector field with $Q_i$
independent of $x_i$.  On the disorder events of
\Cref{prop:stat-posterior-gap},
\begin{equation}\label{eq:exact-hodge-bound}
 \inf_g\E_\mu\sum_i v_i(Q_i-d_ig)^2
 \le C\E_\mu\sum_{i,j}v_iv_j|d_jQ_i-d_iQ_j|^2+o(1).
\end{equation}
Here $C$ depends only on the fixed observation horizon, precision bound,
and posterior gap.  The $o(1)$ is bounded by a polynomial in $n$ times
the exponential stopping probability.  It is taken separately for each
fixed approximant.
\end{lemma}

\begin{proof}
For a full-support posterior $P$ define
\[
 M(P)=\min_g\E_P\sum_i v_{P,i}(Q_i-d_ig)^2,
 \qquad r=Q-dg_P.
\]
Center the minimizer under $P$ to fix its additive constant.  On the
finite cube the Dirichlet form is positive definite modulo constants, so
the minimizer is unique and the following differentiations are legitimate.
We estimate its field derivative, then its observation drift, and finally
the curl energy at the product endpoint.

\emph{The normal equation and its derivative.}
Since $r_i$ is own-spin independent, stationarity on a single edge gives
\begin{equation}\label{eq:exact-normal}
 0=\E_P\sum_i v_{P,i}r_i d_i\phi
   =\E_P\phi\sum_i(x_i-m_{P,i})r_i
 \quad\text{for every }\phi.
\end{equation}
Full support therefore implies
$\sum_i(x_i-m_{P,i})r_i=0$ pointwise.
Freeze the minimizer and vary the external field in coordinate $j$.
The derivative of the normal functional against $\phi$ is exactly
\begin{equation}\label{eq:exact-normal-derivative}
 F_j(\phi)=-\E_P\phi v_{P,j}r_j.
\end{equation}
Indeed the derivative of the measure multiplies the pointwise zero in
\Cref{eq:exact-normal}; the remaining derivative is
$\partial_{h_j}m_{P,i}=\1_{i=j}v_{P,j}$.
Testing \Cref{eq:exact-normal} with $\phi=x_j$ also gives
$\E_Pv_{P,j}r_j=0$.  If the heat-bath gap of $P$ is at least $\gamma$,
then the dual norm relative to $\mathcal E_P$ satisfies
\begin{equation}\label{eq:exact-normal-dual}
 \sum_j\norm{F_j}_{\mathcal E_P^*}^2
 \le\gamma^{-1}\sum_j\E_P(v_{P,j}r_j)^2
 \le\gamma^{-1}M(P).
\end{equation}
This uses Cauchy--Schwarz after centering $\phi$, followed by the
posterior Poincar\'e inequality.

\emph{The drift of the minimum.}
Let $A\succeq0$ be the precision per unit time on one segment of the
observation.  In innovation coordinates the posterior masses satisfy
\[
 \dd P(x)=P(x)(x-\E_Px)^{\mathsf T}A^{1/2}\dd B.
\]
This is the finite-state filtering identity; it follows by the quotient
rule for likelihood weights, as in \cite[Equation~(2.1)]{SKGap}, with $A^{1/2}$ inserted for the stated precision.
For a fixed $i$ and rest configuration $z$, write
$a=P(z,+)$, $b=P(z,-)$, and
\[
 c_i(z)=\frac{4ab}{a+b},\qquad
 v_i(z)=\frac{4ab}{(a+b)^2}.
\]
The edge contribution to the quadratic cost for frozen $r$ is
$c_i(z)r_i(z)^2$.  The Hessian of $c(a,b)=4ab/(a+b)$ is
\[
 D^2c(a,b)=-\frac8{(a+b)^3}
 \begin{pmatrix}b^2&-ab\\-ab&a^2\end{pmatrix}.
\]
The difference of the innovation vectors of the two corners is
$2A^{1/2}e_i$.  It\^o's formula therefore gives precisely
\begin{equation}\label{eq:exact-edge-drift}
 \operatorname{drift}(c_i)=-A_{ii}c_iv_i.
\end{equation}
Summing frozen residual costs bounds this drift below by
$-\norm{A} M(P)$.

The drift of the optimized cost also includes motion of the minimizer.
This contribution is nonpositive and must be bounded from below.  In
coordinates on potentials modulo constants, write the cost near its
minimum as
$M+\mathcal E_P(u)-2F(u)$.
For an infinitesimal innovation, minimization subtracts the squared
dual norm of the first variation $F$.  Thus this contribution is
\[
 -\sum_\alpha\norm{
     \sum_j(A^{1/2})_{j\alpha}F_j
     }_{\mathcal E_P^*}^{\!2}
 \ge-\norm{A}\,\gamma^{-1}M(P),
\]
by \Cref{eq:exact-normal-dual}.  This is also obtained directly by
It\^o differentiating the inverse of the positive Dirichlet matrix.
Together with \Cref{eq:exact-edge-drift}, it proves
\begin{equation}\label{eq:exact-minimum-drift}
 \operatorname{drift}M(P_t)
 \ge-\norm{A_t}(1+\gamma^{-1})M(P_t)
\end{equation}
until the posterior-gap stopping time.

\emph{The product endpoint.}
At a product law write $\sigma_i^2=v_i$ and
$\chi_i=(x_i-m_i)/\sigma_i$.  Expand the own-independent $Q_i$ in the
orthonormal products $\chi_T$.  For each nonempty $S\subset[n]$, put
$b_{i,S}=\sigma_i\widehat Q_i(S\setminus\{i\})$, $i\in S$.
The weighted cost on this support is the squared distance of
$(b_{i,S})_{i\in S}$ from the constant vector, since a potential
coefficient $\widehat g(S)$ contributes that same value to each
$b_{i,S}$.  For $|S|=1$ this distance is zero.  For $|S|\ge2$,
\[
 \sum_{i\in S}|b_{i,S}-\overline b_S|^2
 =\frac1{2|S|}\sum_{i,j\in S}|b_{i,S}-b_{j,S}|^2
 \le\frac14\sum_{i,j\in S}|b_{i,S}-b_{j,S}|^2.
\]
The ordered-pair sum is exactly the weighted curl energy on $S$.
Consequently
\begin{equation}\label{eq:exact-product-hodge}
 M(P_{\rm prod})\le\frac14\E_{P_{\rm prod}}
       \sum_{i,j}v_iv_j|d_jQ_i-d_iQ_j|^2.
\end{equation}

\emph{Transporting the curl energy.}
Fix $i\ne j$ and the remaining spins $z$.  The curl
$c_{ij}(z)=d_jQ_i-d_iQ_j$ is independent of both $x_i,x_j$.
For the four \emph{joint} masses $p_{ab}=P(z,a,b)$ define
\[
 H(p)=\left(\sum_{a,b}p_{ab}^{-1}\right)^{-1}.
\]
This function is homogeneous and concave on the positive cone, since
\begin{equation}\label{eq:exact-harmonic-concavity}
 H(p)=\inf\left\{\sum_{a,b}p_{ab}u_{ab}^2:
                         \sum_{a,b}u_{ab}=1\right\}.
\end{equation}
Extend it continuously to the boundary.  If the pair has no coupling,
$\sum_{a,b}p_{ab}v_i(b)v_j(a)=16H(p)$.
Deleting a coupling $\theta_{ij}$ changes each conditional pair mass
and each conditional variance by factors $e^{O(|\theta_{ij}|)}$,
uniformly in the external fields and $z$.  Hence, throughout the
observation path where $|\theta_{ij}|\le|J_{ij}|$,
\begin{equation}\label{eq:exact-pair-weight}
 e^{-C|J_{ij}|}16H(p)
 \le\sum_{a,b}p_{ab}v_i(b)v_j(a)
 \le e^{C|J_{ij}|}16H(p).
\end{equation}
This concave function of the four joint masses supplies the required
transport estimate.  Each joint posterior mass is a martingale, so
concavity in
\Cref{eq:exact-harmonic-concavity} implies
$\E H(p_t)\le H(p_0)$.  Multiply by $c_{ij}(z)^2$ and sum over $z,i,j$.
Using \Cref{eq:exact-pair-weight} at the initial point and the product
endpoint shows that expected terminal curl energy is at most
$e^{C\max|J_{ij}|}$ times its initial value.

Finally stop before the posterior gap fails, using the test-independent
stopping rule in \Cref{prop:stat-posterior-gap}.  The horizon and precision
norm are fixed.  Equation~\eqref{eq:exact-minimum-drift} makes
$e^{Ct}M(P_t)$ a submartingale up to this stopping time.  Its expected
terminal value is bounded by \Cref{eq:exact-product-hodge} on good paths.
On stopped paths use
$M(P)\le\E_P\sum_i Q_i^2\le\sup_x\sum_iQ_i(x)^2$, a polynomial bound.
Their probability is exponentially small.  More explicitly, let $T$
be the fixed observation horizon, let $\tau$ be the stopping time, and
write
\[
 \mathcal C(P)=\E_P\sum_{i,j}v_{P,i}v_{P,j}
                         |d_jQ_i-d_iQ_j|^2.
\]
Using the full observation posterior $P_T$ also on the exceptional paths,
optional stopping and nonnegativity give
\[
 M(\mu)
 \le e^{CT}\left\{\tfrac14\E\mathcal C(P_T)
                        +n^{C_Q}\P(\tau<T)\right\}
 \le C'\mathcal C(\mu)+o(1).
\]
The first inequality uses the product endpoint on paths that reach $T$
without stopping; the second is the concave curl-energy comparison.
This proves \Cref{eq:exact-hodge-bound}.
All operations can first be localized where posterior masses are bounded
away from zero; bounded finite-cube costs and the displayed inequalities
then remove the localization.
\end{proof}

\begin{proposition}[Closedness of the random visible part]
\label{prop:exact-closedness}
At the comparison times of \Cref{lem:exact-good-times},
\begin{equation}\label{eq:exact-closedness}
 q_s\in\mathcal X_v\mathbin{\widehat\otimes}\mathcal S_s.
\end{equation}
The statement concerns the vector factor of the random-coefficient
projection and requires no regularity of its scalar coefficients.
\end{proposition}

\begin{proof}
Apply the preceding results to a deterministic vector obtained by
contracting $q_s$ with an arbitrary fixed scalar Hilbert direction.
First, \Cref{lem:exact-form-closure} gives strong approximants with
vanishing curl and own-spin derivative.  \Cref{lem:exact-own-spin} makes the approximants own-spin independent,
preserving their weighted limit and vanishing curl.  Apply
\Cref{lem:exact-no-defect} separately to each fixed polynomial approximant,
letting dimension tend to infinity before improving the approximation.
Their weighted distance to exact gradients tends to zero.
\Cref{lem:exact-saturation} puts the limiting deterministic vector in
$\mathcal X_v$.  Orthogonality to $\mathcal X_v^\perp$ after contraction
against every scalar direction gives \Cref{eq:exact-closedness}.
\end{proof}

\subsection{Unweighted decay and the spectrum below the spin edge}

\begin{lemma}[Upgrading decay to the unweighted norm]
\label{lem:exact-upgrade}
Let a $K$-invariant subspace of $\mathcal X_v$ have weighted decay
$\norm{K_tq}_v\le C e^{-at}\norm{q}_v$.  For every
$\rho<\min\{a,1-\beta^2\}$ and every $q\in\mathcal H$ whose weighted
image lies in this subspace,
\begin{equation}\label{eq:exact-unweighted-decay}
 \norm{K_tq}_0\le C_\rho e^{-\rho t}\norm{q}_0.
\end{equation}
The same bound holds with a Hilbert scalar tensor factor and conditional
expectation contractions.  In particular, the full exact space has
unweighted decay at every rate below $\lambda$.
\end{lemma}

\begin{proof}
The weighted estimate alone cannot be inverted through $\iota$.
Instead, separate the gradient extension into its diagonal differential
part
\[
 (\mathcal Dp)_i=(L-I)p_i-2x_i\sum_j a_{ij}d_jp_i
\]
and the multiplication operator $ap$.  Uniformly in $x,i$,
\begin{equation}\label{eq:exact-sharp-row}
 \sum_j\frac{a_{ij}^2}{v_j}
 \le\sum_jJ_{ij}^2(1+o(1))\le\beta^2+o(1).
\end{equation}
The first inequality follows by the one-variable mean-value formula for
$\tanh$, with error uniform because $\max|J_{ij}|\to0$; the second is
uniform concentration of row square sums.  Integration against stationarity
and completion of the square give
\[
 \ip{p}{\mathcal Dp}_0
 \le-\norm{p}_0^2-\mathfrak d(p)
   +2\sqrt{(\beta^2+o(1))\norm{p}_0^2\mathfrak d(p)}
 \le-(1-\beta^2-o(1))\norm{p}_0^2.
\]
Thus its limiting semigroup $T_t^{\rm diag}$ has norm at most
$e^{-(1-\beta^2)t}$.  The other part maps weighted to unweighted norm:
\[
 \norm{ap}_0\le\norm{aD_v^{-1}}\norm{D_vp}_0\le C\norm{p}_v.
\]
Duhamel's formula consequently gives
\[
 \norm{K_tq}_0
 \le e^{-(1-\beta^2)t}\norm{q}_0
 +C\int_0^t e^{-(1-\beta^2)(t-s)}e^{-as}\norm{q}_v\dd s.
\]
Since $\norm{q}_v\le\norm{q}_0$, the convolution gives
\Cref{eq:exact-unweighted-decay}, including equality of the two exponents
by taking strict slack.  Tensoring and conditional expectation do not
increase any norm used here.  Finally
$\lambda\le(1-\beta)^2<1-\beta^2$ for $0<\beta<1$.
\end{proof}

\begin{proposition}[Isolation below the spin edge]
\label{prop:exact-isolation}
The unweighted semigroup satisfies
\begin{equation}\label{eq:exact-essential-rate}
 \limsup_{t\to\infty}\frac1t
       \log\norm{K_t|_{\mathcal X}}_{\rm ess}
 \le-\lambda_{\rm sp}.
\end{equation}
Every weighted spectral point of $\mathcal A_v$ strictly below
$\lambda_{\rm sp}$ is an isolated eigenvalue of finite multiplicity, and
its eigenspace lifts into $\mathcal X$.  If
$\lambda<\lambda_{\rm sp}$, let $E_v$ be the bottom eigenspace and $E$ its
unweighted lift.  The weighted orthogonal projection $\Pi_v$ and the
unweighted spectral projection $P$ agree through the embedding:
\begin{equation}\label{eq:exact-projection-agreement}
 \Pi_v\iota=\iota P.
\end{equation}
The slow projection also defines a bounded map on every $q\in\mathcal H$
with $\iota q\in\mathcal X_v$: project $\iota q$ onto $E_v$ and lift
through the finite-dimensional inverse $E_v\to E$.  On the weighted
complement there is decay with some rate strictly larger than $\lambda$,
and hence unweighted decay at every rate below the minimum of that rate
and $1-\beta^2$, for all such $q$.
\end{proposition}

\begin{proof}
There are two steps.  First, the hidden estimate bounds the essential
norm in the ordinary exact space.  Second, the dense intertwining map
identifies the resulting finite-dimensional spectral subspaces in the
weighted completion.

\emph{Essential norm.}
For every fixed deterministic $u\in\mathcal H$ and bounded recipe $B$,
the finite-dimensional contractions $u^{\mathsf T}B$ concentrate at
$\ip{u}{B}_{\mathcal H}$.  This first holds for a fixed finite recipe;
for a dense approximation $u'$ use
\[
 \norm{(u-u')^{\mathsf T}B}_{L^2(\mu)}
 \le\sup_x\norm{B(x)}\norm{u-u'}_0.
\]
It need not hold for an arbitrary random scalar times a visible vector.

We turn this concentration into the squared orthogonality required by
the hidden theorem.  Let $u_k$ be a weakly null unit sequence in
$\mathcal X$, and approximate each $u_k$ within $1/k$ by a finite exact
test combination.
At dimension $n_k$ impose concentration for its contractions with the
first $k$ recipes, its norm, and the operator pairing currently being
tested.  Take $n_k$ sufficiently large that the finite pointwise bounds
of this approximant are bounded by a fixed power of $n_k$.  A sufficiently
slow diagonal sequence $k=k(n)\to\infty$ then has
\[
 \E_\mu|u_{k(n)}^{\mathsf T}B|^2\longrightarrow0
 \quad\text{for every fixed recipe }B.
\]
Thus the representatives are hidden in the squared-contraction sense of
\Cref{thm:hidden-propagation}; weak convergence by itself was not used
as a substitute for this property.

To see that this controls the essential norm, take finite-rank
orthogonal projections $P_m^0\uparrow I$ on $\mathcal X$.  The Hilbert-space
identity
\[
 \norm{K_T}_{\rm ess}=\lim_m\norm{(I-P_m^0)K_T}
\]
follows because $P_m^0K_T$ is finite rank and every compact operator is
uniformly small after applying $I-P_m^0$.  If the claimed essential-rate
bound failed, choose unit $u_m\in\operatorname{ran}(I-P_m^0)$ and unit
inputs $p_m$ witnessing the excess.  The outputs are weakly null, so the
previous diagonal construction applies to these arbitrary witnesses;
approximate the inputs too by finite exact test combinations while
preserving the pairing.  Split $K_T=K_{T-1}K_1$ for $T>1$.
\Cref{lem:exact-smoothing} supplies a uniform stationary Hessian budget
for the exact inputs $K_1p_m$.  Since their derivatives are independent
of the differentiating spin, their $w_j$-weighted budget equals their
$v_j$-weighted budget.  Apply \Cref{thm:hidden-propagation} to the hidden
outputs and these smoothed inputs.  Here the hidden theorem imports the
multitime conditional derivative argument: its fixed smooth circuits
use the mixed-slot hierarchy, common tensor truncations, and opened
network reduction of
\Cref{lem:path-mixed-tensors,lem:path-conditional-derivative,prop:path-opened-derivative}.
Its path entry comparisons retain the typical integrated-parameter
convention of \Cref{thm:path-calculus}; no supremum over auxiliary
Duhamel times is added here.  The bounded $K_1$ factor and loss of one
time unit disappear after dividing the logarithm by $T$.
This proves \Cref{eq:exact-essential-rate} for all possible norm
witnesses.

\emph{The dense embedding and its spectral projections.}
Fix $0<\rho<\lambda_{\rm sp}$ and choose $T$ so large that the essential
norm of $T_X:=K_T|_{\mathcal X}$ is smaller than $e^{-\rho T}$.
Choose $r>e^{-\rho T}$ away from the spectrum, with $r<1$.
The usual resolvent argument for a finite-rank perturbation of an
operator of norm smaller than $r$ gives a Riesz decomposition
\[
 \mathcal X=E_r\oplus F_r,
\]
where $E_r$ is finite dimensional, contains the spectrum outside $|z|=r$,
and the spectral radius of $T_X|_{F_r}$ is at most $r$.
For clarity, a finite-rank approximation with smaller remainder norm
makes the resolvent outside that norm the inverse of a finite-dimensional
analytic matrix; its zeros are isolated and have finite multiplicity.
On a compact annulus $|z|\ge r$ there are finitely many such zeros.
This supplies precisely the decomposition just used.

Let $T_v=K_T|_{\mathcal X_v}$.  It is positive self-adjoint and
$\iota T_X=T_v\iota$.  Injectivity implies that every generalized
eigenvector in $E_r$ embeds as a generalized eigenvector of $T_v$.
Self-adjointness removes all nontrivial Jordan chains and forces these
eigenvalues to be real and positive.  Write $W_r=\iota E_r$.
For $f\in F_r$ and $r'>r$, the spectral-radius formula yields
\[
 \norm{T_v^k\iota f}_v
 \le\norm{\iota}\,C_{r'}(r')^k\norm{f}_0.
\]
If the $T_v$-spectral measure of $\iota f$ charged $(r',\infty)$, its
$k$th moment in this inequality would grow faster than $(r')^k$.
The spectral theorem therefore forces that measure into $[0,r']$;
let $r'\downarrow r$.  Thus $\iota F_r$ is orthogonal to $W_r$.
On the dense image $\iota\mathcal X$ the Riesz decomposition is already
an orthogonal decomposition into $W_r$ and the low spectral subspace.
It follows that
\[
 \mathcal X_v=W_r\mathbin{\oplus^{\perp}}
                    \overline{\iota F_r},
 \qquad
 \Pi_{W_r}\iota=\iota P_{E_r}.
\]
Density now excludes additional weighted spectral mass outside $[0,r]$.
The displayed identity proves agreement of the projections through
$\iota$; no boundedness of an ordinary Riesz projection in the weaker
norm was assumed.

Given any rate strictly below $\lambda_{\rm sp}$, choose $\rho$ above
that rate and then $r$ between the essential disk and its corresponding
semigroup eigenvalue.  The preceding argument accounts for all weighted
spectrum below that rate by finitely many eigenvectors in $\mathcal X$.
If the bottom $\lambda$ is strictly below $\lambda_{\rm sp}$, it is
therefore attained and isolated.  Its complement has weighted generator
spectrum bounded below by $\lambda+\eta$ for some $\eta>0$.
The inverse $E_v\to E$ is bounded because these spaces are finite
dimensional, so its composition with $\Pi_v\iota$ gives the asserted
projection for vectors of $\mathcal H$ whose image is weighted-exact.
Applying \Cref{lem:exact-upgrade} to their complementary parts proves the
last assertion.
\end{proof}

\begin{remark}
The construction of $\lambda$ uses stationary fixed-time limits alone.\label{rem:exact-parameter-order}
The subsequent use of \Cref{prop:exact-closedness} on a nonstationary
trajectory chooses a fixed interior smoothing interval, a finite set of
recipes, derivative and likelihood truncations, and the required
conditional simulation accuracies before dimension.  Only after those
limits are established are the finite bases enlarged and spectral or
large-block limits taken.  This order permits arbitrary scalar
coefficients in \Cref{eq:exact-transport} and keeps the nonequivalent
Hilbert norms distinct throughout.
\end{remark}
